\chapter{Symbolic Domain Subset Testing\label{chap:SymbolicDomainSubsetTesting}}

Whether an assignment statement is well-typed depends on whether the dynamic domain of the
right-hand side type is contained in the dynamic domain of the left-hand side type,
for any given dynamic environment.

\begin{definition}[Domain Subset]
For any given types $\vt$ and $\vs$ and \staticenvironmentterm{} $\tenv$,
we say that $\vt$ is a \emph{domain subset} of $\vs$ in $\tenv$,
if the following condition holds:
\hypertarget{def-domainsubset}{}
\begin{equation}
\begin{array}{l}
\domainsubset(\tenv, \vt, \vs) \triangleq \\
\qquad \forall \denv\in\dynamicenvs.\
\dynamicdomain((\tenv, \denv), \vt) \subseteq \dynamicdomain((\tenv, \denv), \vs) \enspace.
\end{array}
\end{equation}
\end{definition}

For example, consider the assignment
\begin{center}
\texttt{var x : $\overname{\texttt{integer\{1,2,3\}}}{\vs}$ = ARBITRARY : $\overname{\texttt{integer\{1,2\}}}{\vt}$;}
\end{center}

It is well-typed, since the dynamic domain of \verb|integer{1,2,3}| is\\
$\{\nvint(1), \nvint(2), \nvint(3)\}$ in every dynamic environment, which is a superset of
the dynamic domain of \verb|integer{1,2}|, which is $\{\nvint(1), \nvint(2)\}$ in every dynamic environment.

Since dynamic domains are potentially infinite, this requires \emph{symbolic reasoning}.
Furthermore, since any (\symbolicallyevaluableterm{}) expressions may appear inside integer and bitvector
types, domain subset testing is undecidable.
We therefore approximate domain subset testing \emph{conservatively} via the predicate $\symdomsubsettest(\tenv, \vt, \vs)$.

\hypertarget{def-symdomsubsettest}{}
\begin{definition}[Sound Domain Subset Test]
A predicate
\[
  \symdomsubsettest(\overname{\staticenvs}{\tenv} \aslsep \overname{\ty}{\vt} \aslsep \overname{\ty}{\vs}) \aslto \Bool
\]
is \emph{sound} if the following condition holds:
\begin{equation}
  \begin{array}{l}
  \forall \vt,\vs\in\ty.\ \tenv\in\staticenvs. \\
  \;\;\;\; \symdomsubsettest(\tenv, \vt, \vs) \typearrow \True \;\implies\; \domainsubset(\tenv, \vt, \vs)  \enspace.
  \end{array}
\end{equation}
\end{definition}

That is, if a sound domain subset test returns a positive answer, it means that
$\vt$ is definitely a domain subset of $\vs$ in the \staticenvironmentterm{} $\tenv$.
This is referred to as a \emph{true positive}.
However, a negative answer means one of two things:
\begin{description}
  \item[True Negative:] indeed, $\vt$ is not a domain subset of $\vs$ in the \staticenvironmentterm{} $\tenv$; or
  \item[False Negative:] the symbolic reasoning is unable to decide.
\end{description}

In other words, $\symdomsubsettest(\tenv, \vt, \vs)$ errs on the \emph{safe side} ---
it never answers $\True$ when the real answer is $\False$, which would (undesirably)
determine the following statement as well-typed:
\begin{center}
\texttt{var x : $\overname{\texttt{integer\{1,2\}}}{\vs}$ = ARBITRARY : $\overname{\texttt{integer}}{\vt}$;}
\end{center}

A sound but trivial domain subset test is one that always returns $\False$.
However, that would make all assignments be considered as not well-typed.
Indeed, it has the maximal set of false negatives.
Reducing the set of false negatives requires stronger symbolic reasoning algorithms,
which inevitably leads to higher computational complexity.
%
The symbolic domain subset test in \chapref{SymbolicDomainSubsetTesting}
attempts to accept a large enough set of true positives, based on empirical trial and error,
while maintaining the computational complexity of the symbolic reasoning relatively low.
%
In particular, it serves as the definitive domain subset test that must be utilized
by any implementation of the ASL type system.

This chapter is concerned with implementing a \hyperlink{def-symdomsubsettest}{sound domain subset test}
for \integertypesterm{} and \bitvectortypesterm, as defined above and as employed by \\
\TypingRuleRef{SubtypeSatisfaction}.
This is technically achieved by first transforming types (in the case of \integertypesterm{})
and width expressions (in the case of \bitvectortypesterm{}) into symbolic representations
that we refer to as \emph{symbolic domains} and then checking subsumption over the symbolic domains.

\ChapterOutline
\begin{itemize}
  \item \secref{Symbolic Domains} defines a symbolic representation for sets of integers; and
  \item \secref{Symbolic Reasoning} defines the symbolic reasoning used for domain subset
        testing.
\end{itemize}

\section{Symbolic Domains\label{sec:Symbolic Domains}}
\hypertarget{def-symbolicdomain}{}

We define the \emph{symbolic domain} datatype, reusing $\intconstraint$ from the untyped AST,
as follows:

\RenderTypes{symbolic_domains}


\begin{itemize}
  \item We refer to an element of the form $\Finite(S)$ as a \emph{symbolic finite set integer domain},
        which represents the (non-empty) set of integers $S$;
  \item We refer to an element of the form $\ConstrainedDom(\vc)$ as a \emph{symbolic constrained integer domain},
        which represents the set of integers given by the constraint $\vcs$; and
  \item We refer to an element of the form $\Top$ as a \emph{symbolic unconstrained integer domain},
        which represents the set of all integers.
\end{itemize}

\section{Symbolic Reasoning\label{sec:Symbolic Reasoning}}

The main rule in this section is \TypingRuleRef{SymdomSubsetUnions}, which defines the function
$\symdomsubsetunions$.

Other helper rules are as follows:
\begin{itemize}
  \item \TypingRuleRef{SymdomNormalize}
  \item \TypingRuleRef{SymdomOfType}
  \item \TypingRuleRef{SymdomOfWidthExpr}
  \item \TypingRuleRef{SymdomOfConstraint}
  \item \TypingRuleRef{SymdomEval}
  \item \TypingRuleRef{SymdomSubset}
  \item \TypingRuleRef{ApproxConstraints}
  \item \TypingRuleRef{ApproxConstraint}
  \item \TypingRuleRef{ApproxExprMin}
  \item \TypingRuleRef{ApproxExprMax}
  \item \TypingRuleRef{ApproxBottomTop}
  \item \TypingRuleRef{IntsetToConstraints}
  \item \TypingRuleRef{ApproxExpr}
  \item \TypingRuleRef{ApproxType}
  \item \TypingRuleRef{ConstraintBinop}
  \item \TypingRuleRef{ApplyBinopExtremities}
  \item \TypingRuleRef{PossibleExtremitiesLeft}
  \item \TypingRuleRef{PossibleExtremitiesRight}
  \item \TypingRuleRef{ConstraintPow}
  \item \TypingRuleRef{ConstraintMod}
\end{itemize}

\TypingRuleDef{SymdomSubsetUnions}
\RenderRelation{symdom_subset_unions}

\ExampleDef{Conservatively Testing Subsumption Between Two Lists of Symbolic Domains}
In \listingref{SymdomOfConstraint}, typechecking the assignment \verb|x=y;|
is performed by testing whether each member of
\[
\Subdomains\left(\left[
  \begin{array}{l}
    \Finite(\{2,3,4\}),\\
    \ConstrainedDom(\ConstraintRange(\ELInt{5}, \AbbrevEVar{\vN})),\\
    \ConstrainedDom(\ConstraintExact(\AbbrevEBinop{\ADD}{\AbbrevEVar{\vN}}{\AbbrevEVar{\vN}}))\\
  \end{array}
\right]\right)
\]
is subsumed (that is, represents a subset as defined by $\symdomsubset$) by a member of
\[
\Subdomains\left(\left[
  \begin{array}{l}
    \Finite(\{1,2,3,4,5,6, 9\}),\\
    \ConstrainedDom(\ConstraintRange(\ELInt{5}, \AbbrevEVar{\vN})),\\
    \ConstrainedDom(\ConstraintExact(\AbbrevEBinop{\MUL}{\AbbrevEVar{\vN}}{\ELInt{2}}))\\
  \end{array}
\right]\right) \enspace.
\]
In this example, this is true as the following hold:
\begin{itemize}
  \item $\Finite(\{2,3,4\})$ is subsumed by $\Finite(\{1,2,3,4,5,6, 9\})$,
  \item $\ConstrainedDom(\ConstraintRange(\ELInt{5}, \AbbrevEVar{\vN}))$ (as a constraint for \verb|x|)
        is subsumed by the same element (as a constraint for \verb|y|), and
  \item $\ConstrainedDom(\ConstraintExact(\AbbrevEBinop{\ADD}{\AbbrevEVar{\vN}}{\AbbrevEVar{\vN}}))$
        is subsumed by \\
        $\ConstrainedDom(\ConstraintExact(\AbbrevEBinop{\MUL}{\AbbrevEVar{\vN}}{\ELInt{2}}))$.
\end{itemize}


\RenderProseAndFormally{symdom_subset_unions}

\TypingRuleDef{SymdomNormalize}
\RenderRelation{symdom_normalize}

\ExampleDef{Normalizing Symbolic Domains}
In \listingref{SymdomOfConstraint}, normalizing the \symbolicdomainterm{} for the type\\
\verb|integer{1..6, 5..N, 9, N*2}| (see \ExampleRef{The Symbolic Domain of a Type})
merges $\Finite(\{1,2,3,4,5,6\})$ and $\Finite(\{9\})$, yielding the \symbolicdomainterm{}
\[
\Subdomains\left(\left[
  \begin{array}{l}
    \Finite(\{1,2,3,4,5,6, 9\}),\\
    \ConstrainedDom(\ConstraintRange(\ELInt{5}, \AbbrevEVar{\vN})),\\
    \ConstrainedDom(\ConstraintExact(\AbbrevEBinop{\MUL}{\AbbrevEVar{\vN}}{\ELInt{2}}))\\
  \end{array}
\right]\right) \enspace.
\]


\RenderProseAndFormally{symdom_normalize}

\TypingRuleDef{SymdomOfType}
\RenderRelation{symdom_of_type}

\ExampleDef{The Symbolic Domain of a Type}
In \listingref{SymdomOfConstraint}, the \symbolicdomainterm{} for the type
\verb|integer{1..6, 5..N, 9, N*2}|
is
\[
\Subdomains\left(\left[
  \begin{array}{l}
    \Finite(\{1,2,3,4,5,6\}),\\
    \ConstrainedDom(\ConstraintRange(\ELInt{5}, \AbbrevEVar{\vN})),\\
    \Finite(\{9\}),\\
    \ConstrainedDom(\ConstraintExact(\AbbrevEBinop{\MUL}{\AbbrevEVar{\vN}}{\ELInt{2}}))\\
  \end{array}
\right]\right) \enspace.
\]


\RenderProseAndFormally{symdom_of_type}

\TypingRuleDef{SymdomOfWidthExpr}
\RenderRelation{symdom_of_width_expr}

\ExampleDef{The Symbolic Domain of a Bitwidth Expression}
The symbolic domain of the bitwidth expression \verb|2 * x|
is:
\[
\Subdomains([\AbbrevConstraintExact{\AbbrevEBinop{*}{\ELInt{2}}{\AbbrevEVar{\vx}}}]) \enspace.
\]
\ASLListing{The symbolic domain of a bitwidth expression}{SymdomOfWidthExpr}{\typingtests/TypingRule.SymdomOfWidthExpr.asl}

\RenderProseAndFormally{symdom_of_width_expr}

\TypingRuleDef{SymdomOfConstraint}
\RenderRelation{symdom_of_constraint}

\ExampleDef{The Symbolic Domain of a Constraint}
In \listingref{SymdomOfConstraint}, considering the type annotation for \verb|x|:
\begin{itemize}
  \item the \symbolicdomainterm{} for the constraint \verb|1..6| is $\Finite(\{1,2,3,4,5,6\})$
  \item the \symbolicdomainterm{} for the constraint \verb|9| is $\Finite(\{9\})$,
  \item the \symbolicdomainterm{} for the constraint \verb|N*2| is \\
        $\ConstrainedDom(\ConstraintExact(\AbbrevEBinop{\MUL}{\AbbrevEVar{\vN}}{\ELInt{2}}))$,\\
        since the premise
        $\symdomeval(\tenv, \AbbrevEBinop{\MUL}{\AbbrevEVar{\vN}}{\ELInt{2}})$ yields $\Top$,
  \item the \symbolicdomainterm{} for the constraint \verb|5..N| is \\
        $\ConstrainedDom(\ConstraintRange(\ELInt{5}, \AbbrevEVar{\vN}))$,\\
        since the premise
        $\symdomeval(\tenv, \AbbrevEVar{\vN})$ yields $\Top$.
\end{itemize}

\ASLListing{The symbolic domain of a constraint}{SymdomOfConstraint}{\typingtests/TypingRule.SymdomOfConstraint.asl}


\RenderProseAndFormally{symdom_of_constraint}

\TypingRuleDef{SymdomEval}
\RenderRelation{symdom_eval}

We assume that $\ve$ has been annotated as it is part of the constraint for an integer type,
and therefore applying $\normalize$ to it does not yield a \typingerrorterm{}.

\ExampleDef{Evaluating Expressions to Construct Symbolic Domains}
In \listingref{SymdomOfConstraint}, considering the constraints in the type annotation for
\verb|x|,
we have that both
$\symdomeval(\tenv, \AbbrevEBinop{\MUL}{\AbbrevEVar{\vN}}{\ELInt{2}})$
and $\symdomeval(\tenv, \AbbrevEVar{\vN})$ yield\\$\Top$,
since \verb|N| is a parameter.
On the other hand, $\symdomeval(\tenv, \ELInt{1})$ and $\symdomeval(\tenv, \ELInt{6})$,
yield $1$ and $6$, respectively.


\RenderProseAndFormally{symdom_eval}

\TypingRuleDef{SymdomSubset}
\RenderRelation{symdom_subset}
The test first \emph{overapproximates} $\visone$ by a set of integers $\vsone$.
That is, $\vsone$ represents a superset of the set of numbers represented by $\visone$.
Notice that this can always be done as $\N$ overapproximates any set of integers.
%
Second, the test \emph{underapproximates} $\vistwo$ by a set of integers $\vstwo$.
That is, $\vstwo$ represents a superset of the set of numbers represented by $\visone$.
Notice that this can always be done as the empty set underapproximates any set of integers.
%
The test concludes by checking whether $\vsone$ is a subset of $\vstwo$,
which if true implies that the set of numbers represented by $\visone$ is a subset of the
set of numbers represented by $\vstwo$ (in any environment).
A negative answer on the other hand means that the
test is unable to conclude subsumption.

\ExampleDef{Conservatively Testing Whether One Symbolic Domain is a Subset of Another}
In \listingref{SymdomOfWidthExpr}, typechecking the assignment \verb|x=y;|
involves test whether the \symbolicdomainterm{} for \verb|2..4| is a subset of the
\symbolicdomainterm{} for \verb|1..6|, which $\True$.
Similarly for testing whether the \symbolicdomainterm{} for \verb|N+N|
is a subset of the \symbolicdomainterm{} for \verb|N*2| also yields $\True$.
In contrast, testing whether the \symbolicdomainterm{} for \verb|5..N|
is a subset of the \symbolicdomainterm{} for \verb|N*2| yields $\False$.

\hypertarget{def-approximationdirectionterm}{}
In the following rules, we use the symbol $\Over$ to stand for \emph{overapproximation},
the symbol $\Under$ to stand for \emph{underapproximation},
both defined by the following type: \RenderType{approximation_direction}
We refer such a symbol as an \approximationdirectionterm.
We use \RenderType{ApproximationFailure} to denote failure to approximate a set of integers.
Specifically, $\CannotOverapproximate$ means that a construct
(like an expression or a constraint) cannot be precisely overapproximated,
and $\CannotUnderapproximate$ means that a construct cannot be precisely
underapproximated.


\RenderProseAndFormally{symdom_subset}

\TypingRuleDef{ApproxConstraints}
\RenderRelation{approx_constraints}

\ExampleDef{Approximating Constraint Lists}
The specification in \listingref{ApproxConstraints} is well-typed,
showing how constraint lists are overapproximated.
\ASLListing{Approximating constraints Lists}{ApproxConstraints}{\typingtests/TypingRule.ApproxConstraints.asl}


\RenderProseAndFormally{approx_constraints}
\CodeSubsection{\ApproxConstraintsBinopBegin}{\ApproxConstraintsBinopEnd}{../Typing.ml}

\TypingRuleDef{ApproxConstraint}
\RenderRelation{approx_constraint}

\ExampleDef{Approximating Constraints}
The specification in \listingref{ApproxConstraint} is well-typed,
showing how constraints are overapproximated and underapproximated.
\ASLListing{Approximating constraints}{ApproxConstraint}{\typingtests/TypingRule.ApproxConstraint.asl}


\RenderProseAndFormally{approx_constraint}

\TypingRuleDef{MakeInterval}
\RenderRelation{make_interval}

\ExampleDef{Approximating Intervals}
The following are examples of over- and under-approximating intervals:
\[
\begin{array}{rcl}
\makeinterval(\Over, 1, 3) &\typearrow& \{1, 2, 3\}\\
\makeinterval(\Under, 1, 3) &\typearrow& \{1, 2, 3\}\\
\makeinterval(\Over, 4, 3) &\typearrow& \CannotOverapproximate\\
\makeinterval(\Under, 4, 3) &\typearrow& \CannotUnderapproximate\\
\end{array}
\]


\RenderProseAndFormally{make_interval}

\TypingRuleDef{ApproxExprMin}
\RenderRelation{approx_expr_min}

\ExampleDef{Approximating the Minimal Value of an Expression}
In \listingref{ApproxExprMin}, approximating the minimal value of
\verb|A DIV 2| for typing the declaration of \verb|y|
yields $\CannotOverapproximate$, since \verb|A| is a parameter with no constraints.
Approximating the minimal value of \verb|a| for the declaration of \verb|b|
yields $12$.
\ASLListing{Approximating the minimal value of an expression}{ApproxExprMin}{\typingtests/TypingRule.ApproxExprMin.asl}


\RenderProseAndFormally{approx_expr_min}

\TypingRuleDef{ApproxExprMax}
\RenderRelation{approx_expr_max}

\ExampleDef{Approximating the Maximal Value of an Expression}
In \listingref{ApproxExprMax}, approximating the maximal value of
\verb|2*N| for typing the declaration of \verb|b|
yields $\CannotOverapproximate$, since \verb|N| is a parameter with no constraints.
Approximating the maximal value of \verb|2*a| for the declaration of \verb|b|
yields $10$.
\ASLListing{Approximating the maximal value of an expression}{ApproxExprMax}{\typingtests/TypingRule.ApproxExprMax.asl}


\RenderProseAndFormally{approx_expr_max}

\TypingRuleDef{ApproxBottomTop}
\RenderRelation{approx_bottom_top}

\ExampleDef{Approximation Extremes}
In \listingref{ApproxBottomTop}, approximating the constraint \verb|a..b|
with $\Over$ yields \\
$\CannotOverapproximate$.
\ASLListing{Approximation extremes}{ApproxBottomTop}{\typingtests/TypingRule.ApproxBottomTop.asl}

\RenderProseAndFormally{approx_bottom_top}

\TypingRuleDef{IntsetToConstraints}
\RenderRelation{intset_to_constraints}

\ExampleDef{Transforming Integer Sets to Constraint Lists}
The specification in \listingref{IntsetToConstraints} is well-typed,
showing how an instance of integer sets transformed to constraint lists.
\ASLListing{Transforming Integer Sets to Constraint Lists}{IntsetToConstraints}{\typingtests/TypingRule.IntSetToConstraints.asl}


\RenderProseAndFormally{intset_to_constraints}

\TypingRuleDef{ApproxExpr}
\RenderRelation{approx_expr}

\ExampleDef{Approximating Expressions}
The specification in \listingref{ApproxExpr} is well-typed,
showing how expressions are overapproximated and underapproximated
for the \subtypesatisfiesterm{} check between \verb|bits(M)| and \verb|bits(N)|,
needed to annotate the expression \verb|x as bits(N)|.
\ASLListing{Approximating expressions}{ApproxExpr}{\typingtests/TypingRule.ApproxExpr.asl}


\RenderProseAndFormally{approx_expr}

\TypingRuleDef{ApproxConstraintBinop}
\RenderRelation{approx_constraint_binop}

\ExampleDef{Approximating Binary Expressions Over Constraint Lists}
The specification in \listingref{ApproxConstraintBinop} is well-typed.
Typechecking the specification involves checking that the type of
the expression \verb|(x :: x)| \subtypesatisfiesterm{} \verb|bits(2 * N)|.
This in turn requires:
\begin{itemize}
\item overapproximating \verb|2*N| by applying $\approxconstraintbinop$
      to $\Over$, $\MUL$, and the constraint lists $[2]$, and $[4,8]$,
      which yields $([8, 16], \PrecisionFull)$.

\item underapproximating \verb|N|, which yields the empty set,
      and then applying \\
      $\approxconstraintbinop$
      to $\Under$, $\MUL$, and the constraint lists $[2]$, and $\emptylist$,
      which yields $(\emptylist, \PrecisionFull)$.
\end{itemize}
Since this static subtype satisfaction test fails, the test is deferred to
the dynamic semantics.

\ASLListing{Approximating binary expressions over constraint lists}{ApproxConstraintBinop}{\typingtests/TypingRule.ApproxConstraintBinop.asl}


\RenderProseAndFormally{approx_constraint_binop}

\TypingRuleDef{ApproxType}
\RenderRelation{approx_type}

\ExampleDef{Approximating Types}
The specification in \listingref{ApproxExpr} is well-typed.
In order to test whether \verb|bits(M)| --- the type of \verb|x| --- \subtypesatisfiesterm{}
\verb|bits(N)|, the typechecker overapproximates \verb|M|.
This involves overapproximating the type of \verb|M|, which is \verb|integer{4,8}|.


\RenderProseAndFormally{approx_type}

\TypingRuleDef{ConstraintBinop}
\RenderRelation{constraint_binop}

For examples of the generated constraints
for the modulus operation see \ExampleRef{Generating Modulus Constraints},
for examples of the generated constraints for exponentiation,
see \ExampleRef{Generating Exponentiation Constraints},
and for examples of the generated constraints for all other operations
see \ExampleRef{Generating Constraints for Binary Operations}.


\RenderProseAndFormally{constraint_binop}
\CodeSubsection{\ConstraintBinopBegin}{\ConstraintBinopEnd}{../Typing.ml}

\TypingRuleDef{ApplyBinopExtremities}
\RenderRelation{apply_binop_extremities}

\ExampleDef{Generating Constraints for Binary Operations}
See \ExampleRef{Generating Range-to-Exact Constraints} and
\ExampleRef{Generating Exact-to-Range Constraints} for examples of generating constraints
when one of the constraints is an \exactconstraintterm{} and the other is a \rangeconstraintterm.

The specification in \listingref{ApplyBinopExtremities} shows examples of the constraints
generated (as type annotations on the left-hand-side of local storage declarations)
for cases where a binary operation involves two \rangeconstraintsterm.
\ASLListing{Generating constraints for binary operations}{ApplyBinopExtremities}{\typingtests/TypingRule.ApplyBinopExtremities.asl}

\ProseParagraph
\OneApplies
\begin{itemize}
  \item \AllApplyCase{exact\_exact}
  \begin{itemize}
    \item $\vcone$ is a constraint for the expression $\va$;
    \item $\vctwo$ is a constraint for the expression $\vc$;
    \item define $\newcs$ as the list containing the constraint for the \binopexpressionterm{} $\AbbrevEBinop{\op}{\va}{\vc}$.
  \end{itemize}

  \item \AllApplyCase{range\_exact}
  \begin{itemize}
    \item $\vcone$ is a range constraint for the expressions $\va$ and $\vb$;
    \item $\vctwo$ is a constraint for the expression $\vc$;
    \item applying $\possibleextremitiesleft$ to $\op$, $\va$, and $\vb$ yields $\extpairs$;
    \item define $\newcs$ as the list containing a constraint $\AbbrevConstraintRange{\AbbrevEBinop{\op}{\vap}{\vc}}{\AbbrevEBinop{\op}{\vbp}{\vc}}$
          for each pair of expressions $(\vap, \vbp)$ in $\extpairs$.
  \end{itemize}

  \item \AllApplyCase{exact\_range}
  \begin{itemize}
    \item $\vcone$ is a constraint for the expression $\va$;
    \item $\vctwo$ is a range constraint for the expressions $\vc$ and $\vd$;
    \item applying $\possibleextremitiesright$ to $\op$, $\vc$, and $\vd$ yields $\extpairs$;
    \item define $\newcs$ as the list containing a constraint $\AbbrevConstraintRange{\AbbrevEBinop{\op}{\va}{\vcp}}{\AbbrevEBinop{\op}{\va}{\vdp}}$
          for each pair of expressions $(\vcp, \vdp)$ in $\extpairs$.
  \end{itemize}

  \item \AllApplyCase{range\_range}
  \begin{itemize}
    \item $\vcone$ is a range constraint for the expressions $\va$ and $\vb$;
    \item $\vctwo$ is a range constraint for the expressions $\vc$ and $\vd$;
    \item applying $\possibleextremitiesleft$ to $\op$, $\va$, and $\vb$ yields $\extpairsab$;
    \item applying $\possibleextremitiesright$ to $\op$, $\vc$, and $\vd$ yields $\extpairscd$;
    \item define $\newcs$ as the list containing a constraint $\AbbrevConstraintRange{\AbbrevEBinop{\op}{\vap}{\vcp}}{\AbbrevEBinop{\op}{\vbp}{\vdp}}$
          for each pair of expressions
          $(\vap, \vbp)$ in $\extpairsab$
          and each pair of expressions
          $(\vcp, \vdp)$ in $\extpairscd$.
  \end{itemize}
\end{itemize}

\FormallyParagraph
\begin{mathpar}
\inferrule[exact\_exact]{
  \vcone = \ConstraintExact(\va)\\
  \vctwo = \ConstraintExact(\vc)
}{
  {
    \begin{array}{r}
  \applybinopextremities(\op, \vcone, \vctwo) \typearrow \\
  \overname{[ \ConstraintExact(\AbbrevEBinop{\op}{\va}{\vc}) ]}{\newcs}
    \end{array}
  }
}
\end{mathpar}

\begin{mathpar}
\inferrule[range\_exact]{
  \vcone = \ConstraintRange(\va, \vb)\\
  \vctwo = \ConstraintExact(\vc)\hva\\
  \possibleextremitiesleft(\op, \va, \vb) \typearrow \extpairs\\
  \newcs \eqdef [(\vap, \vbp) \in \extpairs: \AbbrevConstraintRange{\AbbrevEBinop{\op}{\vap}{\vc}}{\AbbrevEBinop{\op}{\vbp}{\vc}}]
}{
  \applybinopextremities(\op, \vcone, \vctwo) \typearrow
  \newcs
}
\end{mathpar}

\begin{mathpar}
\inferrule[exact\_range]{
  \vcone = \ConstraintExact(\va)\\
  \vctwo = \ConstraintRange(\vc, \vd)\hva\\
  \possibleextremitiesright(\op, \vc, \vd) \typearrow \extpairs\\
  \newcs \eqdef [(\vcp, \vdp) \in \extpairs: \AbbrevConstraintRange{\AbbrevEBinop{\op}{\va}{\vcp}}{\AbbrevEBinop{\op}{\va}{\vdp}}]
}{
  \applybinopextremities(\op, \vcone, \vctwo) \typearrow
  \newcs
}
\end{mathpar}

\begin{mathpar}
\inferrule[range\_range]{
  \vcone = \ConstraintRange(\va, \vb)\\
  \vctwo = \ConstraintRange(\vc, \vd)\hva\\
  \possibleextremitiesleft(\op, \va, \vb) \typearrow \extpairsab\\
  \possibleextremitiesright(\op, \vc, \vd) \typearrow \extpairscd\hva\\
  \newcs \eqdef [(\vap, \vbp) \in \extpairsab, (\vcp, \vdp) \in \extpairscd:
  \AbbrevConstraintRange{\AbbrevEBinop{\op}{\vap}{\vcp}}{\AbbrevEBinop{\op}{\vbp}{\vdp}}]
}{
  \applybinopextremities(\op, \vcone, \vctwo) \typearrow
  \newcs
}
\end{mathpar}
%\RenderProseAndFormally{apply_binop_extremities}

\TypingRuleDef{PossibleExtremitiesLeft}
\RenderRelation{possible_extremities_left}

\ExampleDef{Generating Range-to-Exact Constraints}
The specification in \listingref{PossibleExtremitiesLeft}
shows the constraints inferred for various binary operations
(in the type annotations of the corresponding left-hand-side declarations)
where the type of \verb|x| contains a \rangeconstraintterm{} and the type of \verb|c|
contains an \exactconstraintterm.
\ASLListing{Generating range-to-exact constraints}{PossibleExtremitiesLeft}{\typingtests/TypingRule.PossibleExtremitiesLeft.asl}

\RenderProseAndFormally{possible_extremities_left}

\TypingRuleDef{PossibleExtremitiesRight}
\RenderRelation{possible_extremities_right}

\ExampleDef{Generating Exact-to-Range Constraints}
The specification in \listingref{PossibleExtremitiesRight}
shows the constraints inferred for various binary operations
(in the type annotations of the corresponding left-hand-side declarations)
where the type of \verb|a| contains an \exactconstraintterm{} and the type of \verb|y|
contains a \rangeconstraintterm{}.
\ASLListing{Generating exact-to-range constraints}{PossibleExtremitiesRight}{\typingtests/TypingRule.PossibleExtremitiesRight.asl}

\RenderProseAndFormally{possible_extremities_right}

\TypingRuleDef{ConstraintMod}
\RenderRelation{constraint_mod}

\ExampleDef{Generating Modulus Constraints}
The specification in \listingref{ConstraintMod}
shows the constraints inferred for the exponentiation expression \verb|x MOD y|
(in the type annotation for \verb|z|).
\ASLListing{Generating modulus constraints}{ConstraintMod}{\typingtests/TypingRule.ConstraintMod.asl}

\ExampleDef{Ill-typed Modulus Constraint Assignment}
In \listingref{typing-constraintmod}, the assignment to \texttt{z} is illegal, since the type
inferred for \texttt{z} is\\
\verb|integer{0..2}|.

\ASLListing{An ill-typed modulus constraint assignment}{typing-constraintmod}{\typingtests/TypingRule.ConstraintMod.bad.asl}


\RenderProseAndFormally{constraint_mod}
\CodeSubsection{\ConstraintModBegin}{\ConstraintModEnd}{../Typing.ml}

\TypingRuleDef{ConstraintPow}
\RenderRelation{constraint_pow}

\ExampleDef{Generating Exponentiation Constraints}
The specification in \listingref{ConstraintPow}
shows the constraints inferred for the exponentiation expression \verb|x^y|
(in the type annotation for \verb|z|).
\ASLListing{Generating exponentiation constraints}{ConstraintPow}{\typingtests/TypingRule.ConstraintPow.asl}


\RenderProseAndFormally{constraint_pow}
\CodeSubsection{\ConstraintPowBegin}{\ConstraintPowEnd}{../Typing.ml}
